\chapter{Discussion}
\label{ch:discussion}

% ─────────────────────────────────────────────────────────────────────────────
\section{Does the System Answer the Research Question?}
\label{sec:disc_rq}
% ─────────────────────────────────────────────────────────────────────────────

The central research question asks whether a hybrid IRT 3PL + BKT adaptive
engine, implemented in a production web application, can deliver personalised
question sequencing producing measurable within-session learning gains for UK
KS3/GCSE Geography.

The system provides a technically complete answer. The adaptive engine correctly
selects items by Fisher information at the learner's current $\hat{\theta}$,
enforces Bloom-level prerequisites, tracks per-KC BKT knowledge state, and
applies the credit factor modification for hint-assisted responses. All
components behave as specified, confirmed by the algorithm validation evidence
in Track~1 (Section~\ref{sec:results_algorithm}).

The empirical answer is provisionally supported by the within-session
$\Delta\hat{\theta}$ data from $n = 4$ participants, but cannot be stated with
statistical confidence. The honest position is: the system is \textit{technically
validated} (behaves correctly) and \textit{empirically indicative} (preliminary
data is directionally consistent with the expected effect), but \textit{not
yet inferentially confirmed} (a study with $n \geq 15$ and a control condition
is required for a publishable effect size estimate). This distinction between
technical and empirical validation is frequently elided in student project
reports and is made explicit here.

% ─────────────────────────────────────────────────────────────────────────────
\section{Comparison with Related Work}
\label{sec:disc_comparison}
% ─────────────────────────────────────────────────────────────────────────────

\textcite{vanlehn2011} identifies step-based feedback as the critical
determinant of ITS effectiveness ($d = 0.76$), contrasting with answer-based
systems ($d = 0.31$). GeoMentor provides question-level feedback (correct/incorrect
plus explanation) rather than step-level feedback (reasoning step diagnosis),
strictly classifying it as answer-based under VanLehn's taxonomy. Two features
move it toward the step-based category: the three-level hint system
progressively scaffolds the solution pathway before the learner commits an
answer, and the Misconceptions tab identifies the specific wrong-answer
attractor selected, enabling targeted reflection on the specific reasoning
error made. Whether these features produce step-based effect sizes remains an
empirical question. VanLehn's $d = 0.76$ benchmark may be an optimistic
comparator for this system's interaction granularity.

\textcite{kulik2016} report $d = 0.66$ for continuous real-time adaptation
versus $d = 0.28$ for checkpoint systems, directly supporting the design
decision to implement per-response EAP and BKT updating. This is one area
where the implementation is at the upper bound of what the literature predicts
for systems of this class.

% ─────────────────────────────────────────────────────────────────────────────
\section{Scope Reduction: Strength Not Weakness}
\label{sec:disc_scope}
% ─────────────────────────────────────────────────────────────────────────────

The decision to scope down from multi-domain to single-domain Geography
produces more defensible research contributions. \textcite{desmarais2012}
note that the most common failure mode in adaptive system research is shallow
multi-domain validation, in which a system is demonstrated to work across many
domains with weak evidence in each. Deep single-domain validation (rigorous
psychometric calibration, detailed misconception analysis, and multi-track
validity evidence) produces stronger contributions. The geographic focus also
enables a distinct finding that the multi-domain version could not have made:
demonstrating that ITS methods transfer to a declarative-knowledge non-STEM
domain addresses Gap~1 from the literature review. Furthermore, the
architecture is domain-agnostic by design: the adaptive engine contains no
Geography-specific logic, and extending to additional domains requires only
new KC graphs and question banks.

% ─────────────────────────────────────────────────────────────────────────────
\section{Limitations}
\label{sec:disc_limitations}
% ─────────────────────────────────────────────────────────────────────────────

\subsection{Sample Size and Absence of a Control Condition}
\label{subsec:disc_n4}

The most significant limitation is $n = 4$ external participants. Paired
$t$-test at $n = 4$ has approximately 19\% statistical power at $d = 0.76$
and $\alpha = 0.05$, effectively indistinguishable from chance in terms of
hypothesis testing. The multi-method evaluation approach was adopted precisely
because inferential statistics are not viable at this sample size.

The within-subjects control arm (static question-sequencing, items sorted by
ascending difficulty $b$) was implemented in the system but not completed:
as documented in Section~\ref{subsec:disc_condition_gap}, all evaluation
participants received the adaptive condition due to a signup pipeline gap.
A two-arm study with confirmed random assignment is the most important
methodological addition for any follow-up investigation.

\subsection{Condition Assignment Implementation Gap}
\label{subsec:disc_condition_gap}

A late-stage audit revealed that the \texttt{studyCondition} field, added to
the \texttt{User} model at Sprint~8, was never populated by the signup API
route. The Prisma schema default (\texttt{"adaptive"}) silently applied to
every registered user. The bug was identified by asking ``does the
implementation actually do what the report claims?'' and reading the relevant
source file---a form of claim-against-code audit that is underused in student
project development. The fix was applied before final submission; the evaluation
data could not be retracted.

A post-evaluation data point provides directional evidence. After the primary
evaluation window closed and the signup pipeline was corrected, one additional
participant (static condition, $n = 1$) completed a 15-question session,
recording $\Delta\hat{\theta} = -0.080$ at 60\% accuracy. Against the adaptive
cohort mean of $\Delta\hat{\theta} = +0.479$, this single observation is
consistent with the hypothesis that adaptive sequencing produces superior
within-session ability gains ($\Delta = +0.559$ logits in favour of adaptive).
This observation is $n = 1$, post-hoc, and uncontrolled; no inferential
claim is made. It is reported for transparency and as a signal for the
powered study design.

Furthermore, one evaluation participant (P2) returned for a voluntary second
session. His session-opening $\hat{\theta} = -0.209$ matched his session-1
closing estimate rather than the cold-start prior of $-0.780$, confirming that
cross-session BKT state persistence is functioning correctly.

\subsection{Other Limitations}

\textbf{IRT parameters from small calibration sample.} The discrimination
parameter $a$ is sensitive to sparse data: with fewer than 100 responses per
item, standard errors are large. A recalibration using evaluation session data
is the immediate next step for any continued development.

\textbf{Single session per participant.} The BKT cross-session decay detection
and the 7-day delayed retention assessment could not be empirically evaluated.
Multi-session longitudinal data collection is the most important methodological
gap for future work.

\textbf{Convenience sample generalisability.} University of Hull students
have higher prior Geography knowledge, greater familiarity with
computer-mediated learning, and different motivational profiles than secondary
school learners. The evaluation findings should be interpreted as
proof-of-concept validation in a convenience sample.

% ─────────────────────────────────────────────────────────────────────────────
\section{Threats to Validity}
\label{sec:disc_threats}
% ─────────────────────────────────────────────────────────────────────────────

\begin{description}
  \item[Internal validity.] Absence of a control condition: learning gains
    could reflect content exposure, practice effects, or regression to the mean
    rather than the adaptive mechanism. The within-subjects design partially
    mitigates this; without a completed static arm the mitigation is incomplete.
  \item[Construct validity.] EAP $\hat{\theta}$ is a latent construct; its
    validity depends on IRT calibration quality. Convergent validity would be
    strengthened by comparison with an external geography knowledge measure.
  \item[Ecological validity.] Researcher-administered conditions may produce
    demand characteristics; self-directed deployment would provide more
    ecologically valid engagement data.
  \item[Statistical conclusion validity.] With $n = 4$, statistical conclusion
    validity is severely limited. All findings are descriptive; the risk of
    Type II error is high.
\end{description}

% ─────────────────────────────────────────────────────────────────────────────
\section{What Could Not Be Delivered and Why}
\label{sec:disc_not_delivered}
% ─────────────────────────────────────────────────────────────────────────────

\textbf{Multi-domain support.} The original proposal described coverage of
STEM and Humanities domains. Multi-domain calibration would have required
separate IRT parameter estimates per domain, each demanding a minimum of 200
responses per item, infeasible within the available timeline. The architecture
supports multi-domain extension; the content does not yet exist.

\textbf{Completed A/B randomised controlled trial.} The static branch is
implemented but the condition assignment pipeline was not fully wired at signup
(Section~\ref{subsec:disc_condition_gap}). The comparison relies on VanLehn's
meta-analytic benchmark as a surrogate.

\textbf{$n = 15$ participant target.} Contributing factors include: the
compressed timeline following the scope reduction; reliance on university
participant pools with lower-than-anticipated response rates; and the exclusion
criterion (no prior GCSE Geography) eliminating a substantial proportion of
Hull Computer Science students. These shortcomings are disclosed following
\textcite{woolf2010}'s recommendation that ITS evaluation papers should
explicitly document the gap between planned and achieved validation evidence.

% ─────────────────────────────────────────────────────────────────────────────
\section{Future Work}
\label{sec:disc_future}
% ─────────────────────────────────────────────────────────────────────────────

\begin{enumerate}[leftmargin=2em]
  \item \textbf{Powered evaluation study.} Recruit $n \geq 18$ with a
    completed adaptive/static control condition and 7-day delayed retention
    assessment, transforming proof-of-concept findings into publishable evidence.
  \item \textbf{mirt recalibration.} Rerun IRT calibration using evaluation
    session data to produce more stable item parameter estimates.
  \item \textbf{Cross-Domain Knowledge Transfer.} Extend the KC graph to a
    second domain following the KLI framework of \textcite{koedinger2012} and
    the cross-domain transfer model of \textcite{singley1989}.
  \item \textbf{Step-level feedback upgrade.} Redesign feedback from
    post-hoc explanation (answer-based) to guided reasoning disclosure
    (step-based), targeting VanLehn's $d = 0.76$ category.
  \item \textbf{Longitudinal multi-session study.} Design a three-session
    protocol to evaluate the BKT decay detection mechanism and successive
    relearning boost empirically.
  \item \textbf{Secondary school deployment.} Pilot with UK KS3 students
    in partnership with a geography department, providing ecologically valid
    data with the intended population.
\end{enumerate}
